% !TeX root = ../../../main.tex
\section{模型假设}
\begin{enumerate}
  \item 在此填写与题意一致、建模所必需的假设，并给出依据。
  \item 说明简化条件及其可能影响，并在后文安排检验。
\end{enumerate}

\section{符号说明}
\begin{table}[htbp]
  \centering
  \caption{主要符号及含义}\label{tab:symbols}
  \begin{tabularx}{.84\textwidth}{cXc}
    \toprule
    符号 & 含义 & 单位\\
    \midrule
    $x_i$ & 第 $i$ 个观测值 & 依变量定义\\
    $\hat y_i$ & 模型对第 $i$ 个样本的预测值 & 依变量定义\\
    $n$ & 样本数量 & 无量纲\\
    $\lambda$ & 正则化参数 & 无量纲\\
    \bottomrule
  \end{tabularx}
\end{table}

同一符号在全文中应保持含义一致，首次出现时说明其定义。可以使用表~\ref{tab:symbols} 这样的交叉引用，避免手动维护表格编号。
